\begin{abstract}



随着光伏发电等分布式能源快速发展，如何实现分布式能源本地消纳、减少并网压力，并在满足小区负载的同时降低外部电网购电费用，已经成为微网购电与储能调度调控策略的重要问题。



\textbf{针对问题一}，在电价、小区负载及全天光伏发电功率预测已知的条件下，将一天划分为144个10 min时段，建立以\textbf{全天购电费用最小}为目标的\textbf{混合整数线性规划模型（MILP）}，联合优化外部电网购电与储能充放电。求得\textbf{全天计划购电量为59,482.699 kWh，最低购电费用为35,126.949元}，24:00储电量恢复至6000 kWh。

\textbf{针对问题二}，在负载与光伏出力存在预测偏差、供电不足需按交易时刻电价 5 倍紧急购电的条件下，首先利用实际数据构造完全信息基准；随后基于历史数据进行日前预测，通过滚动MAE选取最近两个同星期历史日平均曲线，结合季节性、近期净负荷偏差及日初储电量修正预测并\textbf{动态设置安全冗余}，形成\textbf{Dynamic日前购电策略}。统计期内，Dynamic紧急购电量为\textbf{318,742.269 kWh}，实际总购电费用为\textbf{15,069,856.025元}；相较固定70 kWh冗余策略，紧急购电量\textbf{减少约79.7\%}，总费用\textbf{减少约639.1万元}。

\textbf{针对问题三}，将0:00、6:00、12:00和18:00的光伏预报插值至10 min尺度，于0:00制定全天计划，并利用后续预报对\textbf{尚未执行时段滚动调整}。调整时保持0:00确定的储能充放电方案不变，仅根据费用比较调整剩余购电量。

\textbf{针对问题四}，引入附件4逐日波动电价，并利用历史电价预测目标日各时段价格。对于问题二，比较不同电价预测方法后，采用\textbf{同星期历史日均值法（WD2）}，其预测MAE为\textbf{0.045153元/kWh}，据此重新制定日前购电与储能策略；对于问题三，根据日内新增光伏预报及已获得的实际电价更新剩余时段预测电价，最终采用\textbf{冻结0:00储能计划的bfreeze策略}。相较不调整方案，该策略使统计期紧急购电量\textbf{下降60.75\%}，总购电费用由16,301,694.650元降至\textbf{16,172,629.246元}，据此确定波动电价下的购电策略。


% 摘要正文结束

\vfill

\noindent\textbf{关键词}：混合整数线性规划；储能调度；动态安全冗余；滚动优化；波动电价

\vspace{0.5cm}


\end{abstract}
